%==========================================================
%  CHAPITRE 3 — Environnement de travail
%  CRM Intelligent avec IA pour Centres d'Appels
%==========================================================

\chapter{Environnement de travail}
\chapterplan{%
\chapterplanitem{1}{Introduction}{sec:chap3_introduction}
\chapterplanitem{2}{Architecture Logicielle}{sec:chap3_architecture_logicielle}
\chapterplanitem{3}{Outils et Technologies}{sec:chap3_outils_technologies}
\chapterplanitem{4}{Environnement de développement}{sec:chap3_environnement_developpement}
\chapterplanitem{5}{Conclusion}{sec:chap3_conclusion}
}

%----------------------------------------------------------
\section*{Introduction}
\phantomsection
\label{sec:chap3_introduction}
\addcontentsline{toc}{section}{Introduction}
%----------------------------------------------------------

\noindent Ce chapitre présente l'environnement complet retenu pour la réalisation de la
plateforme CRM intelligente développée dans le cadre du projet 3LM Solution.
Il s'articule autour de trois axes principaux : le choix d'une architecture clean architecture
avec .NET 8, l'arsenal technologique sélectionné (backend, frontend, IA, DevOps), ainsi que
l'environnement matériel et logiciel qui a soutenu le déroulement des différents sprints.

%==========================================================
\section{Architecture Logicielle}
\label{sec:chap3_architecture_logicielle}
%==========================================================

\noindent L'adoption d'une \textbf{architecture en couches} avec le pattern Repository et
Unit of Work constitue le socle architectural de la plateforme. Ce choix garantit une
séparation claire des responsabilités, une testabilité accrue et une maintenabilité à long
terme.

%----------------------------------------------------------
\subsection{Description de l'Architecture}
%----------------------------------------------------------

\noindent L'architecture repose sur les principes de \textbf{Clean Architecture} et de
\textbf{Domain-Driven Design}~: chaque couche possède une responsabilité bien définie et
les dépendances pointent vers l'intérieur.

\begin{itemize}[label=---]
    \item \textbf{Couche Présentation (API)}~: contrôleurs ASP.NET Core exposant des
    endpoints REST. Valide les entrées, orchestre les appels aux services et retourne
    des réponses standardisées.
    \item \textbf{Couche Application (Services)}~: contient la logique métier, les DTOs
    et les profils AutoMapper. Implémente les cas d'utilisation du système.
    \item \textbf{Couche Domaine (Models)}~: entités métier, value objects et interfaces
    de repository. Aucune dépendance externe.
    \item \textbf{Couche Infrastructure (Data)}~: implémentation des repositories avec
    Entity Framework Core, migrations PostgreSQL, DbContext et configuration des entités.
\end{itemize}

%----------------------------------------------------------
\subsection{Avantages de notre choix}
%----------------------------------------------------------

\noindent Le choix de cette architecture est justifié par les avantages suivants :

\begin{itemize}[label=---]
    \item \textbf{Testabilité}~: chaque couche peut être testée indépendamment grâce à
    l'injection de dépendances et aux interfaces.
    \item \textbf{Maintenabilité}~: la séparation des préoccupations facilite la
    compréhension et l'évolution du code.
    \item \textbf{Indépendance technologique}~: le domaine est totalement découplé des
    frameworks (EF Core, ASP.NET), permettant des changements d'infrastructure sans
    impact sur la logique métier.
    \item \textbf{Réutilisabilité}~: les services métier peuvent être réutilisés par
    différents points d'entrée (API REST, jobs, WebSocket).
\end{itemize}

%==========================================================
\section{Outils et Technologies}
\label{sec:chap3_outils_technologies}
%==========================================================

\noindent Le choix de l'arsenal technologique a été guidé par la recherche d'efficacité,
de performance et de cohérence avec l'architecture cible.

%----------------------------------------------------------
\subsection{Développement Backend}
%----------------------------------------------------------

\subsubsection{.NET 8 — ASP.NET Core}

\noindent .NET 8 est le runtime principal du backend de la plateforme. Cette version LTS
(Long-Term Support) offre stabilité, support à long terme et des performances optimisées
\cite{DotNet8}.

\textbf{Avantages~:} ASP.NET Core offre un framework moderne pour la construction d'API
REST, une intégration native avec Entity Framework Core, une injection de dépendances
intégrée et un pipeline middleware hautement configurable \cite{AspNetCore}.

\begin{figure}[H]
    \centering
    \fbox{\parbox{0.4\linewidth}{
    \centering\vspace{0.3cm}
    \textit{[Logo .NET 8]}\\[0.2cm]
    \vspace{0.3cm}
    }}
    \caption{Logo --- .NET 8}
    \label{fig:logo_dotnet}
\end{figure}

\subsubsection{Entity Framework Core}

\noindent Entity Framework Core (EF Core) est l'ORM (Object-Relational Mapping) utilisé
pour la persistance des données. Il permet d'interagir avec PostgreSQL sans écrire de
requêtes SQL \cite{EFCore}.

\textbf{Avantages~:} EF Core offre un mapping objet-relationnel performant, des migrations
automatisées, un chargement eager/lazy configurable et une intégration native avec le
pattern Repository \cite{Npgsql}.

\begin{figure}[H]
    \centering
    \fbox{\parbox{0.4\linewidth}{
    \centering\vspace{0.3cm}
    \textit{[Logo Entity Framework Core]}\\[0.2cm]
    \vspace{0.3cm}
    }}
    \caption{Logo --- Entity Framework Core}
    \label{fig:logo_efcore}
\end{figure}

\subsubsection{AutoMapper}

\noindent AutoMapper est utilisé pour automatiser le mapping entre les entités du domaine
et les DTOs (Data Transfer Objects) \cite{AutoMapper}.

\textbf{Avantages~:} Réduction significative du code de mapping, profils de configuration
centralisés et validation des mappings au démarrage.

\subsubsection{FluentValidation}

\noindent FluentValidation est la bibliothèque de validation utilisée pour les DTOs
d'entrée des contrôleurs \cite{FluentValidation}.

\textbf{Avantages~:} Règles de validation déclaratives, séparation des validations de la
logique métier et intégration automatique avec le pipeline ASP.NET Core.

%----------------------------------------------------------
\subsection{Développement Frontend}
%----------------------------------------------------------

\subsubsection{React 18 et TypeScript}

\noindent React 18 constitue le framework frontend de la plateforme, avec TypeScript pour
le typage statique \cite{React}\cite{TypeScript}.

\textbf{Avantages~:} React offre une architecture composant modulaire, un écosystème riche
et des performances élevées grâce au Virtual DOM. TypeScript apporte la sécurité du typage
et une meilleure maintenabilité du code.

\begin{figure}[H]
    \centering
    \fbox{\parbox{0.4\linewidth}{
    \centering\vspace{0.3cm}
    \textit{[Logo React 18]}\\[0.2cm]
    \vspace{0.3cm}
    }}
    \caption{Logo --- React 18}
    \label{fig:logo_react}
\end{figure}

\subsubsection{Vite et TailwindCSS}

\noindent Vite est l'outil de build utilisé pour le frontend, offrant un rechargement à
chaud instantané \cite{Vite}. TailwindCSS est le framework CSS utilitaire pour le
styling \cite{TailwindCSS}.

\textbf{Avantages~:} Vite offre des temps de build extrêmement rapides comparé aux
bundlers traditionnels. TailwindCSS permet un développement rapide et cohérent avec
des classes utilitaires.

\subsubsection{Radix UI et Recharts}

\noindent Radix UI fournit les primitives d'interface accessibles et non stylisées
(boutons, modales, select, tabs, etc.) \cite{RadixUI}. Recharts est utilisé pour les
graphiques et les tableaux de bord \cite{Recharts}.

\begin{figure}[H]
    \centering
    \fbox{\parbox{0.4\linewidth}{
    \centering\vspace{0.3cm}
    \textit{[Logo Radix UI]}\\[0.2cm]
    \vspace{0.3cm}
    }}
    \caption{Logo --- Radix UI}
    \label{fig:logo_radix}
\end{figure}

%----------------------------------------------------------
\subsection{Base de Données}
%----------------------------------------------------------

\subsubsection{PostgreSQL 16}

\noindent PostgreSQL 16 est le système de gestion de bases de données relationnelles
utilisé pour la persistance de l'ensemble des données métier \cite{PostgreSQL}.

\textbf{Avantages~:} PostgreSQL offre une fiabilité éprouvée, une intégrité transactionnelle
ACID, un support natif des JSON, des index performants et une excellente scalabilité.

\begin{figure}[H]
    \centering
    \fbox{\parbox{0.4\linewidth}{\centering\vspace{0.3cm}\textit{[Logo PostgreSQL 16]}\\[0.2cm]\vspace{0.3cm}}}
    \caption{Logo --- PostgreSQL 16}
    \label{fig:logo_postgresql}
\end{figure}

%----------------------------------------------------------
\subsection{Intelligence Artificielle}
%----------------------------------------------------------

\subsubsection{Ollama et Gemma 3}

\noindent Ollama est la plateforme d'exécution locale de modèles de langage utilisée
pour le scoring automatique des appels \cite{Ollama}. Le modèle \textbf{gemma3:4b} de
Google DeepMind est utilisé pour l'analyse qualitative des transcriptions \cite{Gemma3}.

\textbf{Avantages~:} L'exécution locale garantit la confidentialité des données audio
et des transcriptions, sans dépendance à une API externe ni coût récurrent. Gemma 3
offre des performances remarquables pour sa taille (4B paramètres).

\begin{figure}[H]
    \centering
    \fbox{\parbox{0.4\linewidth}{\centering\vspace{0.3cm}\textit{[Logo Ollama]}\\[0.2cm]\vspace{0.3cm}}}
    \caption{Logo --- Ollama}
    \label{fig:logo_ollama}
\end{figure}

\subsubsection{Whisper (OpenAI)}

\noindent Whisper est le modèle de reconnaissance vocale (ASR) utilisé pour la
transcription automatique des appels audio \cite{Whisper}.

\textbf{Avantages~:} Whisper supporte nativement le français, offre une excellente
précision de transcription et peut être exécuté localement.

\subsubsection{DistilBERT}

\noindent DistilBERT est utilisé pour l'analyse de sentiment des transcriptions
d'appels \cite{DistilBERT}.

\textbf{Avantages~:} Modèle léger et rapide, adapté à l'inférence en temps réel
avec une précision proche de BERT.

\begin{figure}[H]
    \centering
    \fbox{\parbox{0.4\linewidth}{\centering\vspace{0.3cm}\textit{[Logo Intelligence Artificielle]}\\[0.2cm]\vspace{0.3cm}}}
    \caption{Logo --- Intelligence Artificielle}
    \label{fig:logo_ai}
\end{figure}

%----------------------------------------------------------
\subsection{Sécurité et Authentification}
%----------------------------------------------------------

\subsubsection{JWT (JSON Web Token)}

\noindent JWT est le standard d'authentification utilisé par la plateforme \cite{JWTRFC}.
Les tokens sont signés avec une clé secrète configurable et contiennent l'identité de
l'utilisateur ainsi que ses permissions.

\textbf{Avantages~:} JWT permet une authentification sans état (stateless), idéale pour
les architectures REST. Les tokens peuvent être validés par n'importe quel service sans
requête à une base de données.

\subsubsection{Permission-Based Authorization}

\noindent En complément des rôles, un système de permissions granulaires a été implémenté.
Chaque action (créer un lead, évaluer un appel, configurer les règles, etc.) est protégée
par une permission spécifique.

\textbf{Avantages~:} Ce système offre un contrôle d'accès fin et flexible, permettant de
définir précisément ce que chaque utilisateur peut faire dans la plateforme.

%----------------------------------------------------------
\subsection{Communication Temps Réel}
%----------------------------------------------------------

\subsubsection{WebSocket}

\noindent WebSocket est utilisé pour la messagerie instantanée (chat) et les notifications
en temps réel \cite{WebSocket}.

\textbf{Avantages~:} WebSocket permet une communication bidirectionnelle et persistante
entre le client et le serveur, idéale pour le chat et les mises à jour en direct des
tableaux de bord.

%----------------------------------------------------------
\subsection{DevOps et Hébergement}
%----------------------------------------------------------

\subsubsection{Docker et Docker Compose}

\noindent Docker est utilisé pour la conteneurisation de l'ensemble des services de la
plateforme \cite{Docker}. Docker Compose orchestre les trois conteneurs~: backend .NET 8,
frontend Nginx et PostgreSQL 16.

\begin{figure}[H]
    \centering
    \fbox{\parbox{0.4\linewidth}{\centering\vspace{0.3cm}\textit{[Logo Docker]}\\[0.2cm]\vspace{0.3cm}}}
    \caption{Logo --- Docker}
    \label{fig:logo_docker}
\end{figure}

\subsubsection{GitHub Actions (CI/CD)}

\noindent GitHub Actions automatise les pipelines d'intégration continue et de déploiement
continu \cite{GitHubActions}.

\textbf{Avantages~:} La pipeline exécute les tests unitaires, d'intégration et E2E, vérifie
la couverture de code, scanne les dépendances pour les vulnérabilités et déploie
automatiquement.

\begin{figure}[H]
    \centering
    \fbox{\parbox{0.4\linewidth}{\centering\vspace{0.3cm}\textit{[Logo GitHub Actions]}\\[0.2cm]\vspace{0.3cm}}}
    \caption{Logo --- GitHub Actions}
    \label{fig:logo_github}
\end{figure}

\subsubsection{Swagger / OpenAPI}

\noindent Swagger est utilisé pour la documentation automatique des API REST \cite{Swagger}.

\textbf{Avantages~:} Les endpoints sont documentés et testables directement depuis
l'interface Swagger UI, facilitant l'intégration et les tests.

\subsubsection{Nginx}

\noindent Nginx sert les fichiers statiques du frontend en production et agit comme
reverse proxy \cite{Nginx}.

%----------------------------------------------------------
\subsection{Outils de Développement}
%----------------------------------------------------------

\subsubsection{Visual Studio Code}

\noindent Visual Studio Code est l'IDE principal utilisé pour le développement frontend
(React/TypeScript) et la configuration DevOps \cite{VSCode}.

\subsubsection{Postman}

\noindent Postman est utilisé pour les tests manuels des API REST \cite{Postman}.

\subsubsection{Git}

\noindent Git est le système de gestion de versions utilisé tout au long du projet
\cite{Git}.

%==========================================================
\section{Environnement de développement}
\label{sec:chap3_environnement_developpement}
%==========================================================

\noindent Cette section présente les ressources matérielles et logicielles mobilisées
pour le développement, les tests et le déploiement de la plateforme CRM.

%----------------------------------------------------------
\subsection{Environnement matériel}
%----------------------------------------------------------

\noindent Les développements et tests ont été réalisés sur un poste de travail dont les
caractéristiques sont présentées dans le tableau~\ref{tab:env_materiel}.

\begin{table}[H]
    \centering
    \caption{Caractéristiques de l'environnement matériel}
    \label{tab:env_materiel}
    \renewcommand{\arraystretch}{1.3}
    \begin{tabular}{|l|l|}
        \hline
        \textbf{Composant} & \textbf{Spécification} \\
        \hline
        Processeur & Intel Core i7 (12\textsuperscript{e} génération) \\
        \hline
        Mémoire vive (RAM) & 16 Go minimum (32 Go recommandés) \\
        \hline
        Stockage & SSD 512 Go \\
        \hline
        Système d'exploitation & Windows 11 \\
        \hline
    \end{tabular}
\end{table}

%----------------------------------------------------------
\subsection{Environnement logiciel}
%----------------------------------------------------------

\noindent Le tableau~\ref{tab:env_logiciel} récapitule l'ensemble des outils et
technologies logicielles constituant l'environnement de développement.

\begin{table}[H]
    \centering
    \caption{Environnement logiciel de développement}
    \label{tab:env_logiciel}
    \renewcommand{\arraystretch}{1.3}
    \begin{tabular}{|l|l|}
        \hline
        \textbf{Outil / Technologie} & \textbf{Rôle} \\
        \hline
        .NET 8 SDK & Runtime et framework backend \\
        \hline
        ASP.NET Core & Framework d'API REST \\
        \hline
        Entity Framework Core & ORM et persistance PostgreSQL \\
        \hline
        Npgsql & Provider PostgreSQL pour .NET \\
        \hline
        Node.js 20 / npm & Runtime frontend et gestion des dépendances \\
        \hline
        React 18 / TypeScript & Framework frontend \\
        \hline
        Vite 5 & Bundler et serveur de développement \\
        \hline
        TailwindCSS 3/4 & Framework CSS utilitaire \\
        \hline
        Radix UI & Primitives d'interface accessibles \\
        \hline
        Recharts & Bibliothèque de graphiques \\
        \hline
        Visual Studio Code & IDE principal \\
        \hline
        Docker & Conteneurisation des services \\
        \hline
        PostgreSQL 16 / pgAdmin & Base de données et interface d'administration \\
        \hline
        Ollama / Gemma 3:4b & Moteur d'IA local (scoring) \\
        \hline
        Whisper (OpenAI) & Transcription automatique audio \\
        \hline
        Git / GitHub & Gestion de versions et CI/CD \\
        \hline
        Postman / Swagger & Tests et documentation des API \\
        \hline
        xUnit / Vitest / Playwright & Tests unitaires, composants et E2E \\
        \hline
    \end{tabular}
\end{table}

%==========================================================
\section{Conclusion}
\label{sec:chap3_conclusion}
%==========================================================

\noindent Ce chapitre a présenté l'environnement complet retenu pour la conception et le
développement de la plateforme CRM intelligente. L'adoption d'une architecture en couches
avec .NET 8 et EF Core garantit la maintenabilité et la testabilité du code. Le frontend
React/TypeScript avec TailwindCSS et Radix UI offre une expérience utilisateur moderne et
accessible. Le moteur d'IA local (Ollama, Whisper) assure la confidentialité des données
tout en offrant des capacités avancées d'analyse. Le chapitre suivant abordera la réalisation
des deux premiers sprints constituant le socle de la plateforme.
